\chapter{Architektúra (x86/x64) és tápellátás hatása}

Ebben a fejezetben kifejezetten azokat a kérdéseket vizsgálom, amelyeket különösen érdekesnek tartok:

\begin{itemize}
  \item \textbf{x86 vs x64}: melyik gyorsabb, és miért?
  \item \textbf{Töltőn vs akkumulátorról} (PC3): mennyire látszik a power management hatása?
\end{itemize}

\section{x86 vs x64 -- mit vártam előre?}

\begin{itemize}
  \item A \textbf{natív CPU} útvonalnál az x64 előnye reális: több regiszter, modernebb ABI, jobb optimalizálási lehetőség.
  \item A \textbf{managed .NET} esetén szintén várható különbség, mert a JIT és a natív intrinsics (pl. AES-NI) kihasználása 64 biten tipikusan jobb.
  \item Az \textbf{OpenCL} útvonalnál viszont azt vártam, hogy az x86/x64 különbség kicsi lesz, mert a számítás zöme a GPU-n fut,
        és a host oldalon főleg ``csomagolás'' + adatmozgatás történik. Itt inkább a driver/runtime verzió és a platform számít.
\end{itemize}

A mért adatok ezt nagyrészt alá is támasztják.

\section{PC3: töltő hatása (GCM-256, 64 MB)}

A PC3 egy laptop, ezért különösen érdekes volt megnézni, mennyire fogja vissza a rendszer a teljesítményt akkumulátorról.

\subsection{x64 build: töltőn vs akkuról}

\begin{table}[H]
  \centering
  \caption{PC3 x64 -- átlag throughput \mbps (GCM-256, 64 MB).}
  \fittable{\pgfplotstabletypeset[
    col sep=comma,
    columns={Engine,Direction,ChargerOFF,ChargerON,ratio_on_over_off},
    columns/Engine/.style={string type, column name=Engine},
    columns/Direction/.style={string type, column name=Irány},
    columns/ChargerOFF/.style={column name=Akkuról (\mbps), fixed, precision=2},
    columns/ChargerON/.style={column name=Töltőn (\mbps), fixed, precision=2},
    columns/ratio_on_over_off/.style={column name=Arány (ON/OFF), fixed, precision=3},
    every head row/.style={before row=\toprule, after row=\midrule},
    every last row/.style={after row=\bottomrule},
  ]{\GenPath/pc3_power_gcm256_64mb_x64.csv}}
\end{table}

\begin{figure}[H]
  \centering
  \begin{tikzpicture}
    \begin{axis}[
      ybar,
      bar width=7pt,
      width=0.92\textwidth,
      height=6.2cm,
      enlarge x limits=0.06,
      ylabel={Throughput (\mbps)},
      xtick=data,
      x tick label style={rotate=35, anchor=east},
      legend style={at={(0.5,1.02)},anchor=south,legend columns=2},
      grid=both,
    ]
      \addplot table[x expr=\coordindex, y=ChargerOFF, col sep=comma] {\GenPath/pc3_power_gcm256_64mb_x64.csv};
      \addplot table[x expr=\coordindex, y=ChargerON,  col sep=comma] {\GenPath/pc3_power_gcm256_64mb_x64.csv};
      \legend{Akkuról, Töltőn}
      \node[anchor=north west] at (rel axis cs:0.02,0.98) {\scriptsize Sorrend: Engine/Direction párok, mint a táblában.};
    \end{axis}
  \end{tikzpicture}
  \caption{PC3 x64 -- töltő hatása (GCM-256, 64 MB).}
\end{figure}

\subsection{x86 build: töltőn vs akkuról}

\begin{table}[H]
  \centering
  \caption{PC3 x86 -- átlag throughput \mbps (GCM-256, 64 MB).}
  \fittable{\pgfplotstabletypeset[
    col sep=comma,
    columns={Engine,Direction,ChargerOFF,ChargerON,ratio_on_over_off},
    columns/Engine/.style={string type, column name=Engine},
    columns/Direction/.style={string type, column name=Irány},
    columns/ChargerOFF/.style={column name=Akkuról (\mbps), fixed, precision=2},
    columns/ChargerON/.style={column name=Töltőn (\mbps), fixed, precision=2},
    columns/ratio_on_over_off/.style={column name=Arány (ON/OFF), fixed, precision=3},
    every head row/.style={before row=\toprule, after row=\midrule},
    every last row/.style={after row=\bottomrule},
  ]{\GenPath/pc3_power_gcm256_64mb_x86.csv}}
\end{table}

\subsection{Megfigyelések}

A PC3 mérésekből két hatás nagyon tisztán látszik:

\begin{itemize}
  \item \textbf{Töltőn gyorsabb}: mindhárom engine-nél látszik javulás, különösen a CPU-intenzív (Native/Managed) esetekben.
  \item \textbf{OpenCL-nél kisebb a különbség}: itt is nő a throughput töltőn, de kisebb mértékben.
        Ezt úgy értelmezem, hogy az OpenCL útvonalat inkább a GPU + memória/driver overhead dominálja, míg akkuról elsősorban a CPU órajelek esnek vissza.
\end{itemize}

\section{x86 vs x64 összegzés a mérések alapján}

\subsection{PC2: laptopos konfiguráció (integrált GPU)}

A PC2-n -- amely szintén laptop -- az x64 a natív és managed útvonalon látványosan gyorsabb,
OpenCL-nél viszont csak pár százalék a különbség (a GPU végzi a munkát, a host oldali overhead kisebb arányban változik).

\subsection{PC4: erős diszkrét GPU (RTX 4060)}

PC4-en az OpenCL throughput x86 és x64 között gyakorlatilag azonos (tipikusan 1\% körüli eltérés),
miközben a natív CPU baseline x64-en jelentősen gyorsabb.
Ennek következménye, hogy az OpenCL \emph{speedup} (baseline-hoz viszonyítva) x86 buildben nagyobb számnak látszik,
de ez főleg azért van, mert a baseline lassabb.

\subsection{OpenCL és ``deprecated'' függvények hipotézis}

Saját felvetésem szerint az x64 gyorsabb lehet részben azért, mert az OpenCL kód
kompatibilitási okból régebbi függvényeket is használ az x86-os környezet támogatásához. A \emph{mért adatok} alapján:

\begin{itemize}
  \item ahol ugyanaz az OpenCL platform/driver fut (pl. PC4 NVIDIA CUDA OpenCL), ott az x86/x64 OpenCL throughput szinte azonos,
        tehát a kernel futásideje domináns, és a host architektúra alig látszik.
  \item ahol az OpenCL platform eleve gyengébb (pl. integrált GPU), ott is csak kicsi az arch különbség.
\end{itemize}

Ebből én azt a következtetést vonom le, hogy \textbf{ha van is} ilyen deprecated API hatás,
akkor a jelen benchmarkokban \emph{nem ez} a fő tényező -- a driver/runtime és a GPU képességei sokkal jobban számítanak.

\cleardoublepage
